\section{Bewertung der Implementierung}

Die Implementierung des Systems war technisch anspruchsvoll. Wir haben mehrere Technologien eingesetzt, die wir in der Schule zuvor nicht verwendet hatten, darunter Blazor WebAssembly als Frontend-Framework, Azure AI Search als Suchdienst, die Microsoft Graph API für die SharePoint-Integration, die OpenAI Responses API mit schema-gebundenen Prompts für die Dokumentenklassifikation sowie die Infrastruktur rund um Microsoft Entra ID.
Die Umsetzung ist insgesamt gelungen. Die einzelnen Services sind voneinander getrennt, die Kommunikationswege über REST-Schnittstellen sind klar definiert und sicherheitsrelevante Aspekte wurden berücksichtigt. Die benutzerbezogene Zugriffskontrolle in der Suche und die Trennung zwischen Preflight und Commit im Upload-Flow sind Beispiele dafür. Das Ergebnis ist funktional und entspricht den Anforderungen einer Diplomarbeit.
An zwei Stellen haben wir mehr Aufwand betrieben als notwendig. Alle Services wurden in Microsoft Entra ID registriert und entsprechend konfiguriert, obwohl ein vereinfachter Ansatz ausgereicht hätte. Zusätzlich haben wir Azure AI Search als dedizierten Suchdienst eingesetzt, obwohl die Suchfunktion auch mit der bereits vorhandenen PostgreSQL-Datenbank umsetzbar gewesen wäre. Beide Ansätze erschienen uns zum Zeitpunkt der Planung plausibel, weil sie der Vorgehensweise in produktiven Umgebungen entsprechen. Da alle Services intern betrieben werden und nur von einem einzelnen Client aufgerufen werden, hätte eine einfachere Authentifizierung über API-Keys oder eine zentrale Token-Validierung ausgereicht. Für die Suche hätte PostgreSQL mit integrierter Volltextsuche den Anforderungen genügt, da die Dokumentenbasis überschaubar ist und keine komplexen Ranking-Algorithmen oder semantische Suche benötigt wurden.